\chapter{Reference Data and Symbology Services}\label{ch:pl:reference-data-and-symbology-services}

In January 2022 the ticker META belonged to an exchange-traded fund. At the end of that month the fund changed its ticker to
METV, and on 9 June 2022 a large company that had traded as FB for ten years began trading as META, keeping the same identifier
number for its shares. A system that keyed its history by ticker now held, under one symbol, half a year of a fund followed by
the history of a different issuer, and nothing warned it: both were valid prices of a security that really did trade as META on
those days. This chapter builds the service that prevents that (\texttt{firm.refdata}): it keeps every fact every vendor ever
delivered about an instrument, bitemporally; maps every external symbol to a permanent identifier as of a date and as known at a
time; assembles a \omterm{def:pl:reference-data-and-symbology-services:golden}{golden copy} with the provenance of each value; serves corporate actions as events whose meaning can change
before they happen; and keeps the firm's counterparties the same way. A synthetic universe with planted ticker reuse, a vendor
that misses a rename and a mistyped lot size measures what each piece is worth.

\section{What the reference-data service serves}

One Quant Book~1 (chapter~28) defined reference data as everything about an instrument that is not a price, and symbology as the
set of identifiers it is known by; One Quant Book~7 (chapter~4) built a security master of permanent identifiers and ticker
changes for research. The platform needs the same facts as a service, read by every system of the map of chapter~1: the \omterm{def:pl:a-tick-store-on-open-formats:store}{tick
store} to name its partitions, trading to know a lot size and a tick, risk to find an instrument's underlying, post-trade to find
its ISIN and settlement currency.

\begin{definition}[Reference-data service]\label{def:pl:reference-data-and-symbology-services:service}
A \emph{reference-data service}\index{reference-data service} is the \omterm{def:pl:the-platform-map:sor}{system of record} for instruments, legal entities and
corporate actions: it ingests the facts delivered by data vendors and exchanges, keeps every version of each with the time it
was valid and the time it became known, reconciles the sources, and answers every other system's questions about them as of a
date and as known at a time.
\end{definition}

\begin{omfigure}
\begin{tikzpicture}[font=\scriptsize,>=Stealth,box/.style={draw,rounded corners=2pt,fill=omDef!10,minimum height=0.7cm,inner sep=3pt,align=center}]
\node[box,fill=gray!12] (a) at (0,1.3) {vendor A};
\node[box,fill=gray!12] (b) at (0,0.2) {vendor B};
\node[box,fill=gray!12] (c) at (0,-0.9) {corporate-action\\feed};
\node[box,fill=omMeth!15,text width=3.2cm] (st) at (3.8,0.2) {bitemporal facts\\key $\to$ permanent id};
\node[box,fill=omThm!12,text width=2.3cm] (g) at (7.3,0.9) {golden copy\\with provenance};
\node[box,fill=omThm!12,text width=2.3cm] (x) at (7.3,-0.5) {conflict report\\(operations)};
\node[box,text width=2.9cm] (con) at (10.8,0.2) {tick store, trading,\\risk, post-trade,\\research};
\draw[->] (a) -- (st); \draw[->] (b) -- (st); \draw[->] (c) -- (st);
\draw[->] (st) -- (g); \draw[->] (st) -- (x);
\draw[->] (g) -- (con);
\end{tikzpicture}

\omcaption{The \omterm{def:pl:reference-data-and-symbology-services:service}{reference-data service}: vendors' facts and the corporate-action feed are kept bitemporally under permanent
identifiers; a \omterm{def:pl:reference-data-and-symbology-services:golden}{golden copy} is assembled field by field with the source of each value, disagreements go to a conflict report,
and every consumer asks the service rather than a copy of a vendor file.}\label{fig:pl:reference-data-and-symbology-services:arch}
\end{omfigure}

The chapter's universe has 110 listings over 500 trading days: 100 at the start, 10 listed later. Each has a permanent external
key (its ISIN in the synthetic data), a ticker, a name, a lot size and a tick. Twelve listings change ticker; five of their old
tickers are later given to five of the new listings, as META passed from a fund to a company. Twenty listings split, each
announced 10 to 30 days before the split takes effect; one announced split is cancelled.

\section{Vendors, conflicts and the golden copy}

No single source is right about everything. Vendors differ in coverage, speed and care: one receives exchange notices
directly and updates tickers the same day, another is better on names and legal details; each makes mistakes the other does not.
Two vendors deliver the chapter's facts. Vendor A is prompt, but mistypes one listing's lot size (10 instead of 100) for six
days before correcting it. Vendor B learns two ticker changes two days late, and misses one rename entirely until three days
after its old ticker has been given to a new listing: for 146 days it maps that ticker to the listing that no longer uses it.

\begin{definition}[Golden copy]\label{def:pl:reference-data-and-symbology-services:golden}
The \emph{golden copy}\index{golden copy} of an instrument is the record the \omterm{def:pl:reference-data-and-symbology-services:service}{reference-data service} assembles from its sources,
field by field, by stated rules (a precedence of sources per field, validation, manual overrides), with the source of each
value kept as its provenance; it is what every consumer reads, while the vendors' own records are kept beside it.
\end{definition}

\omcode{code/firm/refdata/firm_refdata.py}{81}{97}{The golden copy (the first source in the field's precedence that has a
value, with its name as provenance) and the conflict report (every instrument whose sources disagree).}\label{lst:pl:reference-data-and-symbology-services:golden}

\Cref{tab:pl:reference-data-and-symbology-services:quality} counts, over the 52\,492 instrument-days of the universe and as known
each day, how often each source and the \omterm{def:pl:reference-data-and-symbology-services:golden}{golden copy} are wrong, and how often the service reports a conflict. With vendor A first
for tickers the \omterm{def:pl:reference-data-and-symbology-services:golden}{golden copy} is never wrong about a ticker, and the conflict report shows every day vendor B was wrong except the
three days it had no value at all. The lot size shows the other side: vendor A is first for lots too, so the \omterm{def:pl:reference-data-and-symbology-services:golden}{golden copy} copies
its typo for all six days. Precedence chooses whose errors the firm inherits; the conflict report is what catches them, since
both sources had a value and they differed on each of those days.

\begin{table}[ht]
\centering\small
\begin{tabular}{lrrrr}
\toprule
field & vendor A wrong & vendor B wrong & \omterm{def:pl:reference-data-and-symbology-services:golden}{golden copy} wrong & conflicts reported \\
\midrule
ticker (A first) & 0 & 153 & 0 & 150 \\
lot size (A first) & 6 & 0 & 6 & 6 \\
\bottomrule
\end{tabular}
\caption{Instrument-days, out of 52\,492 and as known each day, on which each source and the \omterm{def:pl:reference-data-and-symbology-services:golden}{golden copy} disagree with the
truth, and on which the service reports a conflict between the two vendors. Data: \texttt{pl\_refdata.golden\_quality}.}\label{tab:pl:reference-data-and-symbology-services:quality}
\end{table}

The lot typo also shows why the service is bitemporal rather than a table that is updated. Vendor A's typo is a fact valid
from day 100, delivered on day 100; its correction is another version of the same fact, still valid from day 100, delivered on
day 106 (\cref{fig:pl:reference-data-and-symbology-services:bitemporal}). Asked about day 103 as known on day 103, the service
answers 10, the lot a trading system actually used that day; asked about day 103 as known on day 106 or later, it answers 100,
the truth. An updated table can give only the second answer, and every order sized with 10 becomes inexplicable.

\begin{omfigure}
\begin{tikzpicture}[font=\scriptsize,>=Stealth,x=0.9cm,y=0.9cm]
\draw[->] (0,0) -- (8.6,0) node[right]{valid date};
\draw[->] (0,0) -- (0,4.4) node[above]{knowledge time};
\foreach \x/\l in {1/96,3/100,5/103,7/106} {\draw (\x,0.08) -- (\x,-0.08) node[below]{\l};}
\foreach \y/\l in {1/96,2/100,3/103,4/106} {\draw (0.08,\y) -- (-0.08,\y) node[left]{\l};}
\fill[omDef!15] (1,1) rectangle (3,4.2);
\node at (2,2.6) {lot 100};
\fill[omThm!20] (3,2) rectangle (8.3,4);
\node at (5.6,2.5) {lot 10 (A's typo)};
\fill[omMeth!20] (3,4) rectangle (8.3,4.2);
\draw[omMeth,thick] (3,4) -- (8.3,4);
\node[anchor=south,omMeth] at (6.6,4.2) {lot 100 (correction)};
\fill[omThm] (5,3) circle (2pt);
\node[anchor=south west] at (5.05,3.0) {(103, 103): 10};
\fill[omMeth] (5,4.1) circle (2pt);
\end{tikzpicture}

\omcaption{Vendor A's lot size for listing 7, bitemporally: horizontally the date a value is valid for, vertically the time it was
known. The typo is valid from day 100 and known from day 100; its correction is valid from the same day and known from day 106.
A question names a point: about day 103 as known on day 103 the answer is the typo, 10, the value systems used that day; as known
from day 106, the corrected 100.}\label{fig:pl:reference-data-and-symbology-services:bitemporal}
\end{omfigure}

\begin{method}[Running a golden copy]\label{met:pl:reference-data-and-symbology-services:golden}
(1) State a precedence of sources for each field, from their measured accuracy and speed on that field. (2) Keep every source's
value, never only the winner. (3) Report every disagreement to an operations queue the same day, with the values and their
sources, and resolve it by evidence (an exchange notice, the issuer). (4) Record a manual override as a source of its own, with
its author and reason. (5) Measure each source against the resolved truth, and revise the precedence.
\end{method}

\section{Symbol resolution through time}

\begin{definition}[Symbol resolution]\label{def:pl:reference-data-and-symbology-services:resolution}
\emph{Symbol resolution}\index{symbol resolution} is the mapping of an external symbol (a ticker, a vendor code) to the permanent
identifier of the instrument that carried it on a given date, as the firm knew it at a given time.
\end{definition}

Both times matter. The valid date answers ``which company was META on 3 January 2022?'' --- a fund. The knowledge time answers
``what did the firm believe on the morning it traded?'' --- which, for a vendor that learned of a change late, may differ from the
truth. \texttt{resolve} takes both (\cref{lst:pl:reference-data-and-symbology-services:store}), and the service stores every fact
in Book~7's bitemporal store under an effective-date index, so that ``the value in effect on a date'' is the last fact valid at or
before it, in the version known at the time asked.

\omcode{code/firm/refdata/firm_refdata.py}{38}{58}{Effective-dated facts in a bitemporal store: the value in effect on a date is
that of the latest valid date at or before it, in the version known at the time of the question.}\label{lst:pl:reference-data-and-symbology-services:store}

\begin{proposition}[Today's mapping is not a history]\label{prop:pl:reference-data-and-symbology-services:today}
Resolving a symbol at every past date with the mapping in effect today returns the instrument that carries the symbol today. It
differs from point-in-time resolution exactly on the dates when the symbol belonged to another instrument or to none: before a
rename to it, after a rename away from it, and before a reuse of it.
\end{proposition}

\begin{proof}
Today's mapping is a constant function of the date; point-in-time resolution follows the symbol's history. They agree on a date
if and only if the symbol's holder on that date is today's holder.
\end{proof}

On the chapter's universe, sampling every fifth day, resolution with the last day's mapping names a different listing than
point-in-time resolution on 384 of 10\,493 ticker-days, 3.7\%. \Cref{fig:pl:reference-data-and-symbology-services:t056} follows
one ticker through its life: it belongs to one listing until day 33, to none until day 176, and then to the listing that reused
it.

\begin{omfigure}
\begin{tikzpicture}
\begin{axis}[width=12cm,height=5cm,xlabel={trading day},ytick={0,1,2},
  yticklabels={none,first holder (listing 56),reuser (listing 100)},ymin=-0.4,ymax=2.5,xmin=0,xmax=500,
  tick label style={font=\small},label style={font=\small},grid=major,grid style={gray!20},
  legend columns=3,legend style={font=\footnotesize,at={(0.5,-0.35)},anchor=north,draw=none,/tikz/every even column/.append style={column sep=8pt}},
  table/col sep=comma]
\addplot[very thick,omDef,const plot] table[x=day,y=pit]{figdata/platforms/06-reference-data-and-symbology-services/t056.csv};
\addplot[thick,omThm,dashed,const plot] table[x=day,y=today]{figdata/platforms/06-reference-data-and-symbology-services/t056.csv};
\addplot[thick,omMeth,dotted,const plot,yshift=-2pt] table[x=day,y=vendor_b]{figdata/platforms/06-reference-data-and-symbology-services/t056.csv};
\legend{point in time (golden copy),today's mapping,vendor B alone}
\end{axis}
\end{tikzpicture}

\omcaption{What ticker T056 resolves to on each day: point in time through the \omterm{def:pl:reference-data-and-symbology-services:golden}{golden copy}, with today's mapping, and in vendor
B's records alone. The first holder renames on day 33 and a new listing takes the ticker on day 176; today's mapping attributes
the whole period to the reuser, and vendor B keeps the ticker on the first holder until day 179. Data:
\texttt{fig\_refdata.py}.}\label{fig:pl:reference-data-and-symbology-services:t056}
\end{omfigure}

\begin{dated}{2026-09}{A ticker that changed hands}\label{dat:pl:reference-data-and-symbology-services:meta}
Meta Platforms announced on 31 May 2022 that its Class A common stock would trade on Nasdaq under the ticker META from before
the open on 9 June 2022, replacing FB, with its CUSIP number unchanged. The META ticker had belonged to the Roundhill Ball
Metaverse ETF, which changed its ticker to METV at the end of January 2022 (press reports, consulted September 2026).
\end{dated}

\begin{remark}[Which identifier to key on]\label{rem:pl:reference-data-and-symbology-services:keys}
Tickers are reused; national and international security numbers (the ISIN of ISO~6166) change with some corporate actions and
are shared by an instrument's listings on several venues; the Financial Instrument Global Identifier is assigned once, never
changed and never reused. The firm keys everything by its own permanent identifier and stores every external one as a dated
attribute, as One Quant Book~7's security master does.
\end{remark}

\section{Corporate actions as a service}

A corporate action is known before it happens, and what is known can change: a split is announced, confirmed, sometimes amended
or cancelled. A service that stores only the final record of each action cannot say what a strategy knew on the day it traded.

\begin{definition}[Corporate-action event, instrument lifecycle event]\label{def:pl:reference-data-and-symbology-services:event}
A \emph{corporate-action event}\index{corporate-action event} is the reference-data record of a corporate action: its kind
(split, dividend, merger, spin-off), its terms, its effective date and its status (announced, confirmed, amended, cancelled), each
change of status a new version with the time it became known. An \emph{instrument lifecycle event}\index{instrument lifecycle event}
is any dated change to an instrument's reference data: listing, rename, change of lot or tick, corporate action, delisting,
expiry.
\end{definition}

\omcode{code/firm/refdata/firm_refdata.py}{103}{128}{Corporate actions as versioned events: the latest version known at the time
of the question decides, and only confirmed splits enter the adjustment factor.}\label{lst:pl:reference-data-and-symbology-services:actions}

In the chapter's corporate-action feed each split is announced, confirmed three days later, or cancelled. The adjustment factor
between two dates is then a question with a knowledge time: for the first split (a four-for-one, announced on day 386, effective
on day 412) the factor over the two years is 1 as known on day 385 or 386, and 4 from day 389; the cancelled split never enters
it. A backtest that asks the factor with the knowledge time of each decision is point in time; one that asks with today's
knowledge is not, for any signal that uses price levels.

\section{Counterparties, accounts and the rest}

Instruments are half of reference data. The other half is the parties the firm deals with.

\begin{definition}[Counterparty master]\label{def:pl:reference-data-and-symbology-services:cpty}
A \emph{counterparty master}\index{counterparty master} is the reference-data record of the legal entities the firm deals with:
for each, its legal entity identifier (One Quant Book~2, chapter~28), names, parent entity, jurisdiction and the agreements and
accounts attached to it, effective-dated and bitemporal like instrument data.
\end{definition}

Its first use is aggregation: limits and exposures are set on an ultimate parent, so the parent chain must be right on the date
of the exposure and as known when the limit was checked. In the chapter's example a fund is acquired by a bank's subsidiary on
day 250 and the firm learns it on day 260: asked on day 255 about day 255, the fund's ultimate parent is itself; asked on day 262,
the holding company. Both answers are correct, and a credit report produced on day 255 must be reproducible with the first one
(chapter~25).

\begin{example}[The basket that was not held]\label{ex:pl:reference-data-and-symbology-services:basket}
A research note chooses 20 tickers on day 0 (five of which will be renamed and later reused, six of which will split) and asks for
the basket's buy-and-hold return over the 500 days. Point in time, with split factors, the answer is the truth: $+23.5\%$. With
today's ticker mapping and unadjusted prices it is $-16.6\%$, an error of 40.1 points: unadjusted splits alone cost 45.8 points
(the right listings, unadjusted, return $-22.3\%$), and today's mapping alone adds 5.7 (split-adjusted, it returns $+29.2\%$, the
reusers' returns in place of the intended listings').
\end{example}

\begin{omfigure}
\begin{tikzpicture}
\begin{axis}[width=11cm,height=5.6cm,xbar,bar width=9pt,xmin=-30,xmax=38,xlabel={buy-and-hold return over 500 days (\%)},
  ytick={0,1,2,3,4},yticklabels from table={figdata/platforms/06-reference-data-and-symbology-services/basket.csv}{method},
  y dir=reverse,tick label style={font=\small},label style={font=\small},grid=major,grid style={gray!20},
  nodes near coords,nodes near coords style={font=\footnotesize,/pgf/number format/fixed,/pgf/number format/precision=1},
  table/col sep=comma]
\addplot[fill=omDef!45,draw=omDef] table[x=ret,y=k]{figdata/platforms/06-reference-data-and-symbology-services/basket.csv};
\end{axis}
\end{tikzpicture}

\omcaption{The return of the same 20-ticker basket computed five ways: the truth, point in time (equal to it), today's mapping with
split-adjusted prices, the right listings with unadjusted prices, and today's mapping with unadjusted prices. Data:
\texttt{pl\_refdata.basket\_backtest}.}\label{fig:pl:reference-data-and-symbology-services:basket}
\end{omfigure}

\section{Tutorial: two vendors and a reused ticker}\label{tut:pl:reference-data-and-symbology-services:tut}

\begin{tutorial}
\textbf{Goal.} Build the service from two imperfect vendors and a corporate-action feed, measure the \omterm{def:pl:reference-data-and-symbology-services:golden}{golden copy}, and show what
point-in-time resolution is worth to a backtest. \textbf{End state:} \cref{tab:pl:reference-data-and-symbology-services:quality},
\cref{fig:pl:reference-data-and-symbology-services:t056,fig:pl:reference-data-and-symbology-services:basket}.
\begin{tutsteps}
\item \textbf{Universe}: \texttt{pl\_refdata.universe()}: listings, renames, reuses, splits and prices.
\item \textbf{Deliver}: \texttt{build()} feeds vendor A's and vendor B's facts and the corporate-action feed into
\texttt{firm.refdata}, with the planted defects.
\item \textbf{\omterm{def:pl:reference-data-and-symbology-services:golden}{Golden copy}}: \texttt{golden\_quality(rd, u)} for the ticker and the lot size; \texttt{rd.conflicts} on the day of
the first reuse.
\item \textbf{Resolution}: \texttt{mapping\_differs} and the life of ticker T056.
\item \textbf{Backtest}: \texttt{basket\_backtest} five ways.
\end{tutsteps}
\textbf{What to change next.} Put vendor B first for tickers and count the \omterm{def:pl:reference-data-and-symbology-services:golden}{golden copy}'s errors; add a manual override that fixes
the lot typo on its second day and check that the override, not vendor A, is the value's provenance.
\end{tutorial}

\section{Build: the reference-data service}\label{bld:pl:reference-data-and-symbology-services:refdata}

\begin{build}
\bfield{Purpose} The \omterm{def:pl:the-platform-map:sor}{system of record} for instruments, corporate actions and counterparties: the \omterm{def:pl:a-tick-store-on-open-formats:store}{tick store} names partitions with
it (chapter~4), the data-quality rules check against it (chapter~7), trade capture and post-trade read it (chapters~21--22).
\bfield{Interface} \texttt{RefData(precedence)} with \texttt{ingest(vendor, known, valid, key, field, value)}, \texttt{pid},
\texttt{value}, \texttt{golden}, \texttt{conflicts}, \texttt{resolve(field, value, valid, known)}, \texttt{add\_action},
\texttt{actions}, \texttt{factor(pid, start, end, known)}; \texttt{CounterpartyMaster} with \texttt{add}, \texttt{name},
\texttt{ultimate\_parent}.
\bfield{Rules} Every fact is kept with its valid date and knowledge time, corrections are new versions, never edits; every
external key maps to one permanent identifier forever; the \omterm{def:pl:reference-data-and-symbology-services:golden}{golden copy} names its source for each value; every disagreement is
reported; only confirmed corporate actions adjust prices, as known at the time asked.
\bfield{Acceptance tests} \texttt{code/firm/refdata/tests/}: permanent identifiers and effective values per vendor; precedence,
provenance and conflicts; resolution through a rename and a reuse, and history asked today; a correction as a version; actions
announced, confirmed and cancelled and the factor as known; a parent acquired and learned later.
\bfield{Stretch} Manual overrides as a source; validation rules (a lot size must be a multiple of a round lot) that quarantine a
value before it reaches the \omterm{def:pl:reference-data-and-symbology-services:golden}{golden copy}; the security master of Book~7 rebuilt from the service's facts.
\end{build}

\begin{omsources}
\item Meta Platforms, ``Meta Platforms, Inc. to Change Ticker Symbol to `META' on June~9'', press release, 31 May 2022; press
reports of the Roundhill Ball Metaverse ETF's ticker change, 2022.
\item Object Management Group, Financial Instrument Global Identifier (FIGI), \texttt{openfigi.com}; ISO~6166 (ISIN), described by
the Association of National Numbering Agencies.
\end{omsources}

\section{Exercises}

\begin{exercise}[$\star$]\label{exo:pl:reference-data-and-symbology-services:1}
Of 10\,493 ticker-days sampled, 384 resolve to a different listing with today's mapping. What share is that, and on which kinds of
dates does it happen?
\end{exercise}

\begin{exercise}[$\star$]\label{exo:pl:reference-data-and-symbology-services:2}
Meta's press release said its CUSIP number would not change with its ticker. Why does that matter to a firm's reference data, and
why is the CUSIP still not the firm's key?
\end{exercise}

\begin{exercise}[$\star$]\label{exo:pl:reference-data-and-symbology-services:3}
The first split of the universe (four for one) is announced on day 386, confirmed on day 389 and effective on day 412. What is the
adjustment factor over the two years as known on days 385, 386, 389 and 412?
\end{exercise}

\begin{exercise}[$\star\star$]\label{exo:pl:reference-data-and-symbology-services:4}
Vendor A is first for lot sizes and mistyped one for six days. What does the \omterm{def:pl:reference-data-and-symbology-services:golden}{golden copy} serve on those days, what catches the
error, and what rule would have stopped it reaching the \omterm{def:pl:reference-data-and-symbology-services:golden}{golden copy}?
\end{exercise}

\begin{exercise}[$\star\star$]\label{exo:pl:reference-data-and-symbology-services:5}
A fund was acquired on day 250 and the firm learned it on day 260. What is its ultimate parent for day 255 as known on day 255, and
as known on day 262? Which answer should a credit report printed on day 255 show when reproduced a month later?
\end{exercise}

\begin{exercise}[$\star\star$]\label{exo:pl:reference-data-and-symbology-services:6}
Explain the two parts of the basket's 40.1-point error of \cref{ex:pl:reference-data-and-symbology-services:basket}, and why they do
not add up.
\end{exercise}

\begin{exercise}[$\star\star\star$]\label{exo:pl:reference-data-and-symbology-services:7}
\emph{Coding.} Put vendor B first in the ticker's precedence and rerun \texttt{golden\_quality}. How many instrument-days is the
\omterm{def:pl:reference-data-and-symbology-services:golden}{golden copy} wrong, and why not 153?
\end{exercise}

\begin{exercise}[$\star\star\star$]\label{exo:pl:reference-data-and-symbology-services:8}
\emph{Find the flaw.} ``Our backtests use prices adjusted with today's corporate-action factors. That is safe: a split does not
change a company's value, so adjusting history for it cannot leak information.''
\end{exercise}

\section{Problem: The Ticker That Changed Hands}

\begin{problem}[{Weekend problem --- two vendors and a basket}]\label{pb:pl:reference-data-and-symbology-services:1}
The chapter's universe of 110 listings over 500 days, its two vendors, the corporate-action feed and the 20-ticker basket.

\medskip\noindent\textbf{Part I --- The facts.}
\begin{enumerate}
\item How many renames, reuses and splits does the universe hold, and how many splits are cancelled?
\item What two times does every fact carry, and why both?
\item What does vendor B get wrong, and for how long in the worst case?
\item What does vendor A get wrong?
\item Why does the service keep both vendors' values rather than only the \omterm{def:pl:reference-data-and-symbology-services:golden}{golden copy}?
\end{enumerate}

\medskip\noindent\textbf{Part II --- The \omterm{def:pl:reference-data-and-symbology-services:golden}{golden copy}.}
\begin{enumerate}[resume]
\item On how many instrument-days is each vendor wrong about the ticker, and the \omterm{def:pl:reference-data-and-symbology-services:golden}{golden copy}?
\item Why is the \omterm{def:pl:reference-data-and-symbology-services:golden}{golden copy} wrong about the lot size on six days?
\item How many conflicts does the service report for each field?
\item Which two listings conflict on day 176, and why?
\item What would you change in the precedence, and on what evidence?
\end{enumerate}

\medskip\noindent\textbf{Part III --- Resolution and actions.}
\begin{enumerate}[resume]
\item On what share of ticker-days does today's mapping name the wrong listing?
\item Trace ticker T056 through the 500 days.
\item What adjustment factor does the first split give, as known before and after its confirmation?
\item What happens to the cancelled split?
\item What is the fund's ultimate parent for day 255, as known on day 255 and on day 262?
\end{enumerate}

\medskip\noindent\textbf{Part IV --- The verdict.}
\begin{enumerate}[resume]
\item State the \emph{named result}: the basket's return point in time and with today's mapping and unadjusted prices, and the
two parts of the error.
\item Which of the two errors would a researcher notice first, and which is more dangerous?
\item What must a backtest ask the \omterm{def:pl:reference-data-and-symbology-services:service}{reference-data service}, and with which times?
\item What should the conflict report's owner do with the 146-day disagreement?
\item In one sentence: why must reference data be bitemporal?
\end{enumerate}
\end{problem}

\section{Interview questions}

\begin{interviewq}[$\star$ \iqroles{developer, researcher}]\label{iq:pl:reference-data-and-symbology-services:1}
Why should a firm never key its data by ticker?
\end{interviewq}

\begin{interviewq}[$\star\star$ \iqroles{developer}]\label{iq:pl:reference-data-and-symbology-services:2}
Two data vendors disagree about an instrument's lot size. How should the \omterm{def:pl:reference-data-and-symbology-services:service}{reference-data service} decide, and what should it keep?
\end{interviewq}

\begin{interviewq}[$\star\star$ \iqroles{researcher}]\label{iq:pl:reference-data-and-symbology-services:3}
How do corporate actions create look-ahead in a backtest, and how do you avoid it?
\end{interviewq}

\begin{interviewq}[$\star\star$ \iqroles{developer, risk}]\label{iq:pl:reference-data-and-symbology-services:4}
Why does a \omterm{def:pl:reference-data-and-symbology-services:cpty}{counterparty master} need to be bitemporal? Give an example involving limits.
\end{interviewq}

\begin{interviewq}[$\star\star\star$ \iqroles{developer}]\label{iq:pl:reference-data-and-symbology-services:5}
Design a \omterm{def:pl:reference-data-and-symbology-services:service}{reference-data service} for a firm trading equities, futures and options in several countries: sources, identifiers,
storage, \omterm{def:pl:reference-data-and-symbology-services:golden}{golden copy}, and how consumers query it.
\end{interviewq}
